% Pursuing Stacks (cotes 134-2 to 134-8): archive cote 134-4, site page 61 (PDF page 62), typescript page 149.
% typescript. Draft compared: ps3.tex:3104-3168.
% AI-transcribed by Claude Opus 5.5 (claude-opus-5-5, adaptive thinking, effort high), Message Batches API, batch ps2 (msgbatch_01NVwr1Yuo8qJf5rykeGBXPK), 2026-11-12; not yet checked by eye.
\scanpage{134-4/61}
\mspage{149}

that the product of two elements in A be O-connected.  T\over{h}{H}e topos is called
\uline{strictly} \uline{totally} O-\uline{con\inserted{n}ected}  if it is locally connected,  and if moreover
every O-connected object admits a section  -  which (as we saw earlier to-%
day) implies \hand{$\mathcal{A}$} is totally O-connected,  as the wording suggests.  It
amounts to the same to demand that every ``non empty'' object have a section
-  and for this it is enough that the elements in A have a section.  This
latter condition is trivially checked for all standart test categories
I've met so far  (they all have a final object,  and there are maps of the
latter into any other object of A).  A noticeable counterexample here is
$(\hand{$\Delta$}^{\mathrm{f}})\hat{\ }$ (semisimplicial face complexes,  \uline{without} \uline{degeneracies}),  where \xout{xx}
the \uline{weak} test category \hand{$\Delta$}$^{\mathrm{f}}$ hasnt got a final object  (\hand{$\Delta$}$_{\mathrm{o}}$ definitely
\margin{X} is'nt !)  and no \hand{$\Delta$}$_{\mathrm{n}}$ in \hand{$\Delta$}$^{\mathrm{f}}$ \deleted{except \hand{$\Delta$}$_{\mathrm{o}}$ only} has got a section.

\uline{Remark}  I wonder,  when \inserted{$\mathcal{A}$} is totally O-connected, and moreover modelizing
i.e. the Lawvere element is aspheric over e ,  if this implies \hand{$\mathcal{A}$} is
totally aspheric,  and that every element  which is ``non empty'' has a
section (i.e. \uline{strict} total O-connectedness).  \xout{At any rate, the}

Next thing is to look at
\[ W \subset \mathrm{F}\unsure{\mathrm{l}}(\hand{$\mathcal{A}$}) \]
the set of weak equivalences (as defined by non commutative cohomology
of topoi),  and see if it is ``geometric''(page 142).  Condition a') is \xout{xxxx}
clearly satisfied,  there remains the condition b),  namely whether
O-connectedness of an object of \hand{$\mathcal{A}$} (\hand{$\mathcal{A}$} supposed totally O-connected) can
be tested,  using ``intervals'' which are \xout{xx}\uline{aspheric over} e .  More specifi-%
cally,  we want to test that two sections of I belong to the same connec-%
ted component,  using for ``joining'' them intervals which are \xout{xx}aspheric
over e .  The natural idea here is to assume the generating objects \hand{in $A$}
to be aspheric over e (which implies \hand{$\mathcal{A}$} is \uline{totally aspheric},  not only
totally O-connected),  and to use these objects (endowed with suitable
sections) as testing intervals.  This goes through smoothly indeed,  if
we assume  moreover strict total zero-connectedness.  Thus:

\uline{Proposition}  Let \hand{$\mathcal{A}$} be a topos which is strictly totally aspheric
(namely totally aspheric,  and every ``non empty'' object has a section)\hand{:}
\margin{\unsure{cf. lemme} p.268} Then the set W $\subset$ F\unsure{l}(\hand{$\mathcal{A}$})  of weak equivalences in \hand{$\mathcal{A}$} is ``geometric'' ,  and
accordingly,  the homotopy structure h$_{\mathrm{W}}$  defined in terms of aspheric
homotopy intervals,  is the same as the canonical homotopy structure h$_{\unsure{\mathrm{M}}}$
defined in terms of merely O-connected homotopy intervals.  \hand{Moreover, \inserted{for} any
set $A \subset \mathrm{Ob}\,\mathcal{A}$ which is generating and whose objects are 0-connected,
the set of \deleted{inter} 0-connected intervals \deleted{in A} $\underline{I}=(I,\delta_0,\delta_1)$ with $I$
in $A$, generate the homotopy structure $h_W$.}

\margin{61}
